We study the Langevin dynamics of spherical mixed $p$-spin glasses at fixed temperature, on time scales exponential in the dimension $N$. Below the static transition temperature, the dynamics started from the uniform distribution on the sphere stays at total variation distance close to one from the Gibbs measure up to a time exponential in $N$, with probability tending to one. We prove this for pure models of degree at least four at every such temperature, and for mixtures of minimum degree at least four that satisfy an explicit condition on a function of one variable. It extends to initial distributions that are independent of the disorder or depend weakly on it. Between the dynamical and static transition temperatures, we show for general finite mixtures that, with probability exponentially close to one, the generator has exponentially many exponentially small eigenvalues; a fixed part of the centered spin lies in the corresponding slow modes; the stationary spin correlation stays above every level below $q_*^2$ up to an exponential time, where $q_*$ is an explicit root of the plateau equation of the physics literature; the dynamics started from most points of the Gibbs measure stays far from equilibrium; and the two endpoints of the stationary dynamics share mutual information at least linear in $N$. With probability tending to one, the correlation stays above every level below $q_*$. Near the dynamical transition, exponentially many eigenvalues stay below $e^{-cN}$, where $c$ is proportional to an explicit power of the distance to the transition.
